\documentclass[11pt,a4paper]{article}
\usepackage[T1]{fontenc}
\usepackage[utf8]{inputenc}
\usepackage{lmodern}
\usepackage[margin=27mm,headheight=14pt]{geometry}
\usepackage{amsmath,amssymb,amsthm,mathtools}
\usepackage{microtype,booktabs,array,longtable,enumitem,needspace}
\usepackage[dvipsnames]{xcolor}
\usepackage{xurl}
\usepackage[colorlinks=true,linkcolor=MidnightBlue,citecolor=MidnightBlue,urlcolor=MidnightBlue]{hyperref}
\usepackage{fancyhdr}
\pagestyle{fancy}\fancyhf{}
\fancyhead[L]{\small Marginal critical transseries}
\fancyhead[R]{\small Lambert cores and action budgets}
\fancyfoot[C]{\thepage}
\renewcommand{\headrulewidth}{0.3pt}
\setlength{\parskip}{4pt}
\setlength{\parindent}{1.2em}
\setlist{itemsep=3pt,topsep=5pt}
\numberwithin{equation}{section}
\newtheorem{theorem}{Theorem}[section]
\newtheorem{proposition}[theorem]{Proposition}
\newtheorem{lemma}[theorem]{Lemma}
\newtheorem{corollary}[theorem]{Corollary}
\theoremstyle{definition}
\newtheorem{definition}[theorem]{Definition}
\newtheorem{example}[theorem]{Example}
\theoremstyle{remark}
\newtheorem{remark}[theorem]{Remark}
\DeclareMathOperator{\Li}{Li}
\DeclareMathOperator{\Var}{Var}
\newcommand{\N}{\mathbb N}
\newcommand{\R}{\mathbb R}
\newcommand{\C}{\mathbb C}
\newcommand{\E}{\mathbb E}
\newcommand{\Prob}{\mathbb P}
\newcommand{\dd}{\,\mathrm d}
\newcommand{\bs}{\beta_\star}
\newcommand{\Uc}{U_\star}
\newcommand{\ph}{\varphi}
\newcommand{\cJ}{\mathcal J}
\newcommand{\repo}{\texttt{ProveIt}}
\hypersetup{pdftitle={Marginal Critical Transseries: Convergent Lambert-W Charts, Universal Cutoff Corrections, and Certified Action Budgets},pdfauthor={Research draft prepared for Vladimir Reshetnikov},pdfsubject={Critical countable exponential feedback with an exact inverse-cubic action tail}}

\begin{document}
\begin{center}
{\small\scshape Transseries / critical inversion / countable exponential feedback}\par
\vspace{12pt}
{\LARGE\bfseries Marginal Critical Transseries}\par
\vspace{8pt}
{\Large Convergent Lambert--\(W\) Charts,\par Universal Cutoff Corrections,\par and Certified Action Budgets}\par
\vspace{14pt}
{\normalsize Research draft prepared for Vladimir Reshetnikov}\par
{\small 29 September 2026}\par
\end{center}
\vspace{5pt}
\begin{abstract}
A countable exponential-action family with an exact \(j^{-3}\) tail reaches a
critical boundary that is neither an analytic fold nor a nonintegral stable
power singularity. We prove a convergent inverse chart around a Lambert--\(W_{-1}\)
core, with rational coefficient functions, controlled support, and geometric
remainders. We then obtain an all-orders Gaussian--logarithmic expansion of
the original exponential-sector coefficients in a moving coupling window.
The principal cutoff theorem separates two scales: fluctuations have size
\(\sqrt{n\log n}\), but retaining the \(n\)th sector to fixed relative accuracy
requires only \(\sqrt n\) actions. At a comparable cutoff the exact limiting
retained fraction is \(e^{-1/(2\ell^2)}\). Its first two-scale correction is
independent of every finite change in the action weights, after normalization.
At critical coupling this yields a two-term asymptotic formula for the minimal
action budget. An explicit finite-\(n\) inequality, based on positive coefficient
recurrences and an elementary Fourier majorant, supplies a separate genuine
error certificate. Finally, a convergent reduced logarithmic chart for the finite-action
fold explains how those folds approach the marginal boundary. Conventional
proofs, exact finite checks, numerical diagnostics, and further research targets
are provided.
\end{abstract}
\noindent\textbf{Research status.}
The results concern a precisely specified positive, eventually exact power-tail
class. The classical mechanisms---Lagrange inversion, Lambert inversion,
normal attraction, and Poisson extremes---are not claimed as new. The proposed
contribution is the combined explicit marginal theory, particularly the
prefix-independent cutoff correction and its connection to a finite error
certificate. Global publication priority has not been established. No new Lean
formalization or independently reviewed proof is claimed.

\clearpage
{\small\setlength{\parskip}{0pt}\tableofcontents}
\clearpage

\section{Research gap, contribution, and provenance}
\label{sec:intro}

\subsection{The question left by finite-action and stable-action models}
The inspected \repo\ material contains an extensive inversion calculus,
a finite-action moving-fold analysis, and a classification of countable
exponential feedback \cite{repo-index,repo-fold,repo-regularity}.
These address different phenomena. Analytic finite-action sums can produce a
square-root critical inverse. A countable sum may instead reach its own
nonanalytic boundary before an analytic critical-point theorem applies.
The companion draft \emph{Beyond Finite-Action Folds} treats the exact tail
\(a_j\asymp j^{-1-\alpha}\) with \(1<\alpha<2\), including a stable-scale cutoff
law \cite{prior-stable}. The endpoint \(\alpha=2\) is not obtained by simply
substituting into that nonintegral expansion: a pole in the polylogarithm
formula cancels an analytic term and leaves a logarithm.

This article addresses four concrete questions for that marginal endpoint.
Does its inverse still have a convergent finite-generator chart? What is the
large-order law for the original exponential sectors? Does preserving a
coefficient require a cutoff at the fluctuation scale? Can the asymptotic
answer be accompanied by a finite, computable error inequality? The answers
below are affirmative, explicit, negative, and affirmative, respectively.
These are scoped research questions motivated by the inspected corpus, not
claims to have settled a named longstanding conjecture in general transseries
theory.

\subsection{The main distinction: fluctuations versus the largest action}
For the model introduced in Section~\ref{sec:model}, put
\[
 N_n=\sqrt{An},\qquad B_n=\sqrt{AnH_n},\qquad H_n=\log B_n.
\]
Then \(H_n\sim\tfrac12\log n\), so \(N_n/B_n\to0\). The leading results are
\begin{align}
 [q^n]U_{\bs}(q)
 &\sim \frac{\rho^{-n}}{\sqrt{2\pi}\,nB_n},\label{eq:intro-coeff}\\
 \frac{[q^n]U_{\bs}^{[M]}(q)}{[q^n]U_{\bs}(q)}
 &\longrightarrow \exp\!\left(-\frac{1}{2\ell^2}\right)
 \quad\text{if }\frac{M}{N_n}\to\ell\in(0,\infty).\label{eq:intro-cut}
\end{align}
The first scale describes the centered sum of actions. The second describes
its largest individual actions. At this marginal endpoint the two no longer
coincide. In particular, \(M\gg B_n\) is sufficient but unnecessarily costly;
\(M\gg N_n\) is the exact condition for an asymptotically lossless coefficient
truncation.

The first correction is also explicit. With \(\ell_n=M/N_n\) in a compact
subset of \((0,\infty)\),
\begin{equation}
\frac{[q^n]U_{\bs}^{[M]}}{[q^n]U_{\bs}}
=e^{-1/(2\ell_n^2)}
\left\{1+
\frac{\log(2H_n/\ell_n^2)+1-\gamma-2/\ell_n^2}{4H_n}
+O(H_n^{-3/2})\right\}.
\label{eq:intro-correction}
\end{equation}
No finite-prefix constant occurs in this correction. A more general
coupling--cutoff version is proved in Theorem~\ref{thm:cutoff-window}.

\subsection{Relation to established work and limits of novelty}
The substitution reducing our equation to \(z=q\exp(\beta F(z))\) is a standard
Lagrange-inversion construction. Its probability interpretation belongs to
the established connection between coefficient extraction, random allocation,
and conditioned branching structures; Janson gives a broad account, including
infinite-variance cases and extremes \cite{janson}. The distinction between
soft and hard truncation of power tails is also classical
\cite{chakrabarty-samorodnitsky}. Lambert--\(W\) expansions and polylogarithm
continuation are standard special-function material \cite{dlmf-polylog,dlmf-w}.

Accordingly, the leading normal-attraction and extreme-value mechanisms are
not presented as discoveries. What is developed here is a self-contained,
coefficient-level endpoint calculation with convergent local charts,
uniform correction terms, cancellation of finite-prefix data, a minimal-budget
expansion, and an explicit finite error bound. Establishing whether each
combined formula has a previously published equivalent requires a broader
literature review than the targeted search performed for this draft.

\subsection{Source audit}
The repository crosswalk was made against the group
\begin{center}\small
\path{Analysis/Transseries/docs/series-and-transseries/}
\end{center}
at the observed revision
\begin{center}\small
\texttt{3f2d57344dc6f0b3dfa376f9490bfe13d1beb81e}.
\end{center}
The group README, the transseries project README, the countable-feedback
README, and the descriptions of the moving-fold and adjacent research
packages were inspected. The complete approximately seven-hundred-page
canonical volume was not audited. The abstract, scope, and research questions
of the private stable-tail companion were also inspected. That companion is
used to delimit overlap, not as an unproved premise. Every theorem needed
below is established here. No repository file was modified.

\section{Model, critical geometry, and exact coefficient identities}
\label{sec:model}

\subsection{Assumptions and notation}
Throughout, assume
\begin{equation}
 a_j=a j^{-3}+\delta_j\ge0\quad(j\ge1),\qquad
 a>0,\quad a_1>0,\quad \delta_j=0\quad(j>J_0),
\label{eq:weights}
\end{equation}
for some finite \(J_0\). The finite perturbations \(\delta_j\) may have either
sign, provided that the actual weights remain nonnegative. Define
\begin{align}
 F(z)&=\sum_{j\ge1}a_jz^j=a\Li_3(z)+P(z),
 &P(z)&=\sum_{j\le J_0}\delta_jz^j,\label{eq:F}\\
 m&=\sum_{j\ge1}ja_j,
 &\bs&=m^{-1},\label{eq:critical-beta}\\
 A&=\bs a,
 &b&=\frac32+\frac1a\sum_{j\le J_0}j^2\delta_j,\label{eq:Ab}\\
 \nu&=b+\gamma-\frac32,
 &\rho&=e^{-\bs F(1)},\qquad \Uc=\bs F(1).\label{eq:nu-rho}
\end{align}
The second moment \(\sum j^2a_j\) diverges logarithmically, while \(F(1)\) and
\(m\) are finite. We study the small solution of
\begin{equation}
 U_\beta(q)=\beta\sum_{j\ge1}a_jq^j e^{jU_\beta(q)},\qquad U_\beta(0)=0,
 \label{eq:feedback}
\end{equation}
and its finite-action version \(U_\beta^{[M]}\), obtained by replacing \(F\) by
\(F_M(z)=\sum_{j\le M}a_jz^j\). Write
\(u_n(\beta)=[q^n]U_\beta\) and \(u_n^{[M]}(\beta)=[q^n]U_\beta^{[M]}\).

Substituting \(q=e^{-X}\) turns the convergent small-\(q\) expansion into an
integer-action exponential transseries in \(X\). The critical inverse studied
later uses a different variable, \(\epsilon=\log(\rho/q)\). The distinction
between these two expansions is essential.

\begin{proposition}[Exact reduction and coefficient identity]
\label{prop:lagrange}
For every \(\beta>0\), equation~\eqref{eq:feedback} has a unique formal solution
with zero constant term, and an analytic solution near zero. Set
\(z=q e^{U_\beta(q)}\). Then
\begin{equation}
 U_\beta=\beta F(z),\qquad q=z e^{-\beta F(z)},\qquad
 u_n(\beta)=\frac1n[z^n]e^{n\beta F(z)}\quad(n\ge1).
 \label{eq:lagrange}
\end{equation}
The same statements hold with \(F_M\). All coefficients are nonnegative,
\(u_n^{[M]}\) increases with \(M\), and
\(u_n^{[M]}=u_n\) whenever \(M\ge n\).
\end{proposition}
\begin{proof}
Each coefficient of the right side of~\eqref{eq:feedback} depends only on
strictly earlier coefficients of \(U\), proving formal existence and
uniqueness recursively. Analytic existence follows from the implicit function
theorem at \((q,U)=(0,0)\). The substitution gives
\(z=q\exp(\beta F(z))\). Lagrange inversion yields
\[
 [q^n]\beta F(z(q))=\frac\beta n[z^{n-1}]F'(z)e^{n\beta F(z)}.
\]
Differentiating the exponential shows that the right side equals
\(n^{-1}[z^n]e^{n\beta F(z)}\). Positivity, monotonicity, and agreement through
order \(M\) follow directly from this exponential coefficient formula.
\end{proof}

\begin{proposition}[Poisson representation]
\label{prop:poisson}
Let \(Z_j\) be independent Poisson variables of means \(n\beta a_j\), and set
\(S_n=\sum_{j\ge1}jZ_j\). This sum is finite almost surely. Then
\begin{equation}
 u_n(\beta)=\frac{e^{n\beta F(1)}}n\Prob(S_n=n).
 \label{eq:poisson}
\end{equation}
If \(S_{n,M}=\sum_{j\le M}jZ_j\), and
\(\Lambda_{n,M}=n\beta\sum_{j>M}a_j\), then
\begin{equation}
 \frac{u_n^{[M]}(\beta)}{u_n(\beta)}
 =e^{-\Lambda_{n,M}}\frac{\Prob(S_{n,M}=n)}{\Prob(S_n=n)}.
 \label{eq:poisson-ratio}
\end{equation}
\end{proposition}
\begin{proof}
The total mean number of Poisson events is \(n\beta F(1)<\infty\), so only
finitely many events occur almost surely. The probability generating function
of their total weight is \(\exp(n\beta(F(z)-F(1)))\). Extract its \(n\)th
coefficient and use Proposition~\ref{prop:lagrange}. Independence between the
retained and omitted events gives~\eqref{eq:poisson-ratio}.
\end{proof}

\begin{proposition}[Critical boundary]
\label{prop:boundary}
At \(\beta=\bs\), the real map \(q(z)=z e^{-\bs F(z)}\) is strictly increasing
on \((0,1)\) and maps that interval onto \((0,\rho)\). The series for
\(U_{\bs}\) has radius of convergence exactly \(\rho\), and its continuous
boundary value is \(\Uc\).
\end{proposition}
\begin{proof}
For \(0<z<1\),
\[
 q'(z)=e^{-\bs F(z)}\bigl(1-\bs zF'(z)\bigr)>0,
\]
because \(zF'(z)<m\). The endpoint value is \(\rho\). The Poisson identity
bounds \(u_n(\bs)\le\rho^{-n}/n\), giving radius at least \(\rho\).
The logarithmic inverse in Theorem~\ref{thm:lambert-chart} has unbounded first
derivative at this endpoint, so the real branch cannot extend analytically
through \(\rho\). Thus the radius is exactly \(\rho\). Continuity gives the
stated boundary value.
\end{proof}

\section{A convergent Lambert chart at the marginal boundary}
\label{sec:chart}

\subsection{The logarithmic normal form}
Put \(z=e^{-\tau}\), \(\tau>0\), and
\(\epsilon=\log(\rho/q)\). The exact critical equation becomes
\begin{equation}
 \epsilon=\tau+\bs\bigl(F(e^{-\tau})-F(1)\bigr),\qquad
 U_{\bs}(q)=\Uc+\epsilon-\tau.
 \label{eq:epsilon}
\end{equation}

\begin{lemma}[Exact analytic remainder]
\label{lem:normal}
On the principal logarithmic branch near \(\tau=0\),
\begin{equation}
 \epsilon=\frac A2\tau^2\bigl(\log(1/\tau)+b\bigr)+\tau^3 R(\tau),
 \label{eq:normal}
\end{equation}
where \(R\) is holomorphic on \(|\tau|<2\pi\). More explicitly,
\begin{equation}
 \tau^3R(\tau)=\bs\sum_{k\ge3}
 \frac{(-1)^k}{k!}
 \left(a\zeta(3-k)+\sum_{j\le J_0}\delta_jj^k\right)\tau^k.
 \label{eq:R-series}
\end{equation}
For the unperturbed tail, this is
\(A(\tau^3/12-\tau^4/288+\tau^6/86400+\cdots)\).
\end{lemma}
\begin{proof}
The standard integer polylogarithm expansion is
\[
 \Li_3(e^{-\tau})=\zeta(3)-\zeta(2)\tau+
 \frac{\tau^2}{2}\left(\frac32-\log\tau\right)
 +\sum_{k\ge3}\frac{(-1)^k\zeta(3-k)}{k!}\tau^k.
\]
For completeness, subtract the displayed constant, linear, and logarithmic
terms. The third derivative of the remainder is
\[
 -\frac1{e^\tau-1}+\frac1\tau,
\]
which is holomorphic for \(|\tau|<2\pi\), including zero. Its Taylor series
integrates three times to the indicated remainder. The constant and linear
terms in~\eqref{eq:epsilon} cancel because \(\bs m=1\); the finite polynomial
\(P(e^{-\tau})\) supplies the stated quadratic and higher coefficients.
This also proves convergence without a formal limiting argument in the
polylogarithm order. See \cite{dlmf-polylog} for the classical continuation
formulas.
\end{proof}

Let \(c=A/2\). Define the exact core by
\begin{equation}
 w=-W_{-1}\!\left(-\frac{2\epsilon}{ce^{2b}}\right),\qquad
 t=e^{b-w/2}=\sqrt{\frac{2\epsilon}{cw}},\qquad r=w^{-1}.
 \label{eq:core}
\end{equation}
For sufficiently small positive \(\epsilon\), \(w>1\), \(t\to0\), and
\(r\to0\). The core solves
\(\epsilon=ct^2\log(e^b/t)\). The negative real branch \(W_{-1}\), not
\(W_0\), is required to obtain a root tending to zero.

\begin{theorem}[Convergent marginal inverse chart]
\label{thm:lambert-chart}
There is a function \(v(t,r)\), holomorphic near \((0,0)\), such that
\begin{equation}
 \tau=t\bigl(1+v(t,r)\bigr),\qquad
 v(t,0)=v(0,r)=0.
 \label{eq:chart}
\end{equation}
Evaluating \((t,r)\) from~\eqref{eq:core} gives the unique small positive
solution of~\eqref{eq:normal}. Its expansion
\begin{equation}
 v(t,r)=\sum_{j\ge1}V_j(r)t^j
 \label{eq:V}
\end{equation}
converges normally on smaller polydiscs. Each \(V_j\) is a rational function
of \(r\), with a possible pole only at \(r=1\), of order at most \(2j-1\).
Its coefficients depend on only the first \(j\) Taylor coefficients of
\(R/c\). In particular,
\[
 V_1(r)=-\frac{R(0)}c\frac r{1-r}.
\]
There is no convergent Puiseux representation of \(\tau\) in any finite
ramification \(\epsilon^{1/d}\).
\end{theorem}
\begin{proof}
Set \(h=1+v\). Subtract the exact core equation from~\eqref{eq:normal}, divide
by \(ct^2w/2\), and obtain
\begin{equation}
 h^2-1-2rh^2\log h+\frac{2r}{c}t h^3R(th)=0.
 \label{eq:reduced}
\end{equation}
This is holomorphic near \((t,r,h)=(0,0,1)\), and its derivative with respect
to \(h\) there is \(2\). The analytic implicit function theorem gives the
unique chart. At \(t=0\) or \(r=0\), the nearby solution is \(h=1\), proving
the two vanishing statements.

For the coefficient assertion, write \(R(\tau)/c=\sum_{k\ge0}d_k\tau^k\).
At coefficient \(t^j\), the unknown \(V_j\) occurs linearly with coefficient
\(2(1-r)\). Every other term is a polynomial in \(r,d_0,\ldots,d_{j-1}\)
and earlier \(V_i\). In a product of \(k\ge2\) earlier coefficients whose
indices sum to \(j\), the inductive denominator order is at most \(2j-k\).
Division by \(1-r\) therefore gives at most \(2j-1\). Terms carrying an
explicit positive power of \(t\) have no larger denominator order. This proves
the rationality and pole bound by induction.

For positive \(\tau\) sufficiently small, the derivative of the right side
of~\eqref{eq:normal} is
\[
 c\tau\bigl(2\log(1/\tau)+2b-1\bigr)+O(\tau^2)>0.
\]
Hence the chart agrees with the unique real inverse. Moreover,
\begin{equation}
 \tau\sim \frac{2\sqrt\epsilon}{\sqrt{A\log(1/\epsilon)}}.
 \label{eq:leading-inverse}
\end{equation}
Any nonzero convergent Puiseux series has a leading term
\(C\epsilon^\alpha\), \(C\ne0\). The logarithmic growth ratio in
\eqref{eq:leading-inverse} forces \(\alpha=1/2\), whereas
\(\tau/\sqrt\epsilon\to0\), a contradiction.
\end{proof}

\subsection{An explicit holomorphic majorant}
The existence assertion can be made quantitative without locating a complex
critical value. Choose \(R_0<2\pi\), and suppose
\[
 |R(z)|\le K_0,\qquad |R'(z)|\le K_1\quad(|z|\le R_0).
\]
Choose \(T>0\) such that
\begin{equation}
 \frac54T<R_0,\qquad \frac{TK_0}{c}\le\frac18,
 \qquad \frac{T^2K_1}{c}\le\frac18.
 \label{eq:T-conditions}
\end{equation}

\begin{proposition}[Geometric remainder certificate]
\label{prop:chart-bound}
On \(|t|\le T\), \(|r|\le1/8\), the chart may be chosen with \(|v|\le1/4\).
For fixed \(|r|\le1/8\) and \(|t|<T\), truncating~\eqref{eq:V} after \(N\)
terms gives
\begin{equation}
 \left|\tau-t\left(1+\sum_{j=1}^N V_j(r)t^j\right)\right|
 \le \frac{|t|}{4}\frac{(|t|/T)^{N+1}}{1-|t|/T}.
 \label{eq:chart-tail}
\end{equation}
The Taylor expansion in both \(t\) and \(r\) has support in
\(\{(j,k):j,k\ge1\}\) for \(v\).
\end{proposition}
\begin{proof}
Equation~\eqref{eq:reduced} is equivalent to the fixed-point equation
\[
 v=-\frac{v^2}{2}+r(1+v)^2\log(1+v)
   -\frac{rt}{c}(1+v)^3R(t(1+v)).
\]
For \(|v|\le1/4\), \(|\log(1+v)|\le1/3\). Under
\eqref{eq:T-conditions}, the right side has modulus less than \(1/4\).
Its derivative with respect to \(v\) is bounded by
\[
 \frac14+\frac18\frac54\frac53
 +\frac18\left(\frac{3(5/4)^2}{8}+\frac{(5/4)^3}{8}\right)
 <\frac58.
\]
Thus it is a uniform contraction on the closed disc. Iteration produces a
holomorphic fixed point, and Cauchy's estimate in \(t\) gives
\(|V_j(r)|\le\tfrac14T^{-j}\). Summing the geometric tail proves the bound.
Vanishing on both coordinate axes gives the support assertion.
\end{proof}

\begin{example}[The pure trilogarithm]
For \(a_j=j^{-3}\), the first terms, expressed in the exact \(w\)-coordinate,
are
\begin{align}
 \tau={}&t-\frac{t^2}{6(w-1)}
 +\frac{w^2+8w-5}{144(w-1)^3}t^3\notag\\
 &-\frac{9w^3+24w^2-31w+10}{1296(w-1)^5}t^4+O(t^5).
 \label{eq:pure-inverse}
\end{align}
The remainder is uniform when \(r=1/w\) stays in the polydisc above.
If instead \(a_1\) is increased by \(1/2\), the cubic coefficient of
\(\tau^3R(\tau)\) vanishes exactly. Thus the first correction to the exact
Lambert core disappears. Positivity does not prohibit such local
cancellations; the general chart does not assume \(R(0)\ne0\).
\end{example}

\subsection{A finite logarithmic generating grid}
The exact core can itself be expanded convergently in familiar logarithmic
generators. Set
\begin{equation}
 L=\log\frac{Ae^{2b}}{4\epsilon},\qquad
 s=\frac{2\sqrt\epsilon}{\sqrt{AL}},\qquad
 p=\frac{\log L}{L},\qquad u=\frac1L.
 \label{eq:log-generators}
\end{equation}

\begin{corollary}[Convergent three-generator representation]
There is a holomorphic function \(\mathcal A(s,p,u)\) near zero with
\(\mathcal A(0,0,0)=1\) such that
\begin{equation}
 \tau=s\mathcal A(s,p,u).
 \label{eq:three-grid}
\end{equation}
Consequently the critical inverse has an absolutely convergent expansion in
three explicit small generators, although it has no finite Puiseux chart.
Its first logarithmic block is
\begin{equation}
 \tau=s\left[1-\frac{\log L}{2L}
 +\frac{3(\log L)^2-4\log L}{8L^2}
 +O\!\left(\frac{(\log L)^3}{L^3}\right)+O\!\left(\frac{s}{L}\right)\right].
 \label{eq:log-block}
\end{equation}
\end{corollary}
\begin{proof}
The core satisfies \(w-\log w=L\). Write \(w=L(1+Z)\); then
\(Z=p+u\log(1+Z)\). The implicit function theorem gives a holomorphic
\(Z(p,u)\) with \(Z(0,0)=0\). Now
\[
 t=s(1+Z)^{-1/2},\qquad r=\frac{u}{1+Z},
\]
and substitution in~\eqref{eq:chart} proves~\eqref{eq:three-grid}.
Expanding \(Z=p+up+\cdots\) gives~\eqref{eq:log-block}.
\end{proof}

The support claim here is deliberately local and explicit. It concerns a
convergent analytic series in the stated generators, not a construction of
the full field of arbitrary transseries. Conversely, impossibility of a
pure-power Puiseux representation is not an obstruction to this logarithmic
grid.

\section{All-orders Gaussian--logarithmic sector asymptotics}
\label{sec:sectors}

\subsection{The marginal normal scale and coupling window}
For all sufficiently large \(n\), define the large solution
\begin{equation}
 H_n=-\frac12W_{-1}\!\left(-\frac{2}{An}\right),\qquad
 B_n=e^{H_n},\qquad N_n=\sqrt{An}.
 \label{eq:normal-scale}
\end{equation}
Thus \(B_n^2=AnH_n\), \(H_n=\log B_n>1/2\), and
\(B_n=N_n\sqrt{H_n}\). The choice of the large branch is part of the
normalization. Below, \(H,B,N\) abbreviate these three quantities.
For a real parameter \(\lambda\), set
\begin{equation}
 \beta_n=\bs\left(1+\lambda\frac{B_n}{n}\right),\qquad
 \rho_n=e^{-\beta_n F(1)}.
 \label{eq:coupling-window}
\end{equation}
We always restrict \(\lambda\) to a fixed compact set, so \(\beta_n>0\)
eventually. Importantly, \(\rho_n\) is a boundary normalization; for
\(\beta_n>\bs\) it is not asserted to be the radius of convergence of
\(U_{\beta_n}\). The probability identity is valid regardless of that issue.

Use the principal logarithm and put
\begin{equation}
 \Psi_b(t)=\frac{t^2}{2}\bigl(\log(-it)-b\bigr),\qquad \Psi_b(0)=0,
 \qquad \ph(\lambda)=\frac{e^{-\lambda^2/2}}{\sqrt{2\pi}}.
 \label{eq:Psi}
\end{equation}
For real nonzero \(t\), \(\log(-it)=\log|t|-i\pi\operatorname{sgn}(t)/2\).
Define the absolutely convergent integrals
\begin{equation}
 G_k^{(b)}(\lambda)=\frac{1}{2\pi k!}
 \int_{\R}e^{i\lambda t-t^2/2}\Psi_b(t)^k\,\dd t
 \quad(k\ge0).
 \label{eq:Gk}
\end{equation}
Conjugate symmetry makes them real, and \(G_0^{(b)}=\ph\).

\begin{theorem}[Uniform all-orders sector expansion]
\label{thm:allorders}
For each fixed integer \(K\ge0\) and compact real \(\lambda\)-interval,
\begin{equation}
 nB_n\rho_n^n u_n(\beta_n)
 =\sum_{k=0}^{K}\frac{G_k^{(b)}(\lambda)}{H_n^k}
 +O(H_n^{-K-1}),
 \label{eq:allorders}
\end{equation}
uniformly on that interval. The constants may depend on the fixed weights,
\(K\), and the interval. This is an asymptotic expansion; convergence of the
infinite series in \(1/H_n\) is not part of the assertion.
\end{theorem}

We prove the Fourier estimates rather than assuming a local limit theorem.
This also supplies the estimates required for the moving cutoff.

\begin{lemma}[Fourier damping from the exact tail]
\label{lem:damping}
Let
\[
 D(\theta)=\sum_{j\ge1}a_j(1-\cos j\theta),\qquad
 D_M(\theta)=\sum_{j\le M}a_j(1-\cos j\theta).
\]
There are positive constants \(c_1,c_2\) and \(0<\theta_0<1\) such that
\[
 D(\theta)\ge c_1\theta^2\log(1/|\theta|)
 \quad(0<|\theta|\le\theta_0),\qquad
 D(\theta)\ge c_2\quad(\theta_0\le|\theta|\le\pi).
\]
If \(M/N_n\) lies in a fixed compact subset of \((0,\infty)\), the following
bounds hold, with possibly smaller constants, both for \(D\) and \(D_M\):
\begin{align*}
 n\beta_n D(t/B)&\ge c_1t^2 &&(|t|\le\sqrt B),\\
 n\beta_n D(t/B)&\ge c_1B &&(\sqrt B\le|t|\le\theta_0B),\\
 n\beta_n D(t/B)&\ge c_2n &&(\theta_0B\le|t|\le\pi B).
\end{align*}
The displayed statements for \(D_M\) mean that \(D\) is replaced by \(D_M\).
\end{lemma}
\begin{proof}
For \(|j\theta|\le1\), \(1-\cos(j\theta)\ge c(j\theta)^2\).
Summing the exact tail for \(J_0<j\le1/|\theta|\) gives the logarithm.
For a cutoff, sum instead up to \(\min(M,1/|\theta|)\). Bounded initial
terms can be absorbed by taking \(\theta_0\) small; the bound needed is
\(c\theta^2\log\min(M,1/|\theta|)\) when that minimum is large.
For \(|t|\le\sqrt B\), both \(M\) and \(B/|t|\) are at least a fixed
multiple of \(\sqrt B\), eventually. Their minimum has logarithm at least
\(H/3\), and \(n/B^2=1/(AH)\). For
\(\sqrt B\le|t|\le\theta_0 B\), we have \(1/|\theta|\le\sqrt B<M\).
Choose \(\theta_0\) small enough that
\(\theta^2\log(1/\theta)\) increases on \((0,\theta_0]\). Its value at
\(B^{-1/2}\), after multiplication by \(n\), is of order \(B\).
Finally, \(a_1(1-\cos\theta)\) is bounded below on
\([\theta_0,\pi]\). All constants remain uniform when \(\beta_n\) remains
in a compact positive neighborhood of \(\bs\).
\end{proof}

\begin{proof}[Proof of Theorem~\ref{thm:allorders}]
Fourier inversion and Proposition~\ref{prop:poisson} give
\begin{equation}
 nB\rho_n^n u_n(\beta_n)=\frac{1}{2\pi}
 \int_{-\pi B}^{\pi B}
 \exp\left\{n\beta_n\bigl(F(e^{it/B})-F(1)\bigr)-\frac{int}{B}\right\}
 \,\dd t.
 \label{eq:fourier-full}
\end{equation}
The boundary values of the expansion in Lemma~\ref{lem:normal} give, near
zero,
\[
 \bs\bigl(F(e^{-s})-F(1)\bigr)
 =-s+\frac A2s^2\bigl(\log(1/s)+b\bigr)+s^3R(s).
\]
The identity remains valid on either imaginary side, with the specified
logarithm. Substitute \(s=-it/B\). The linear term, after centering at
\(n\), is exactly \(i\lambda t\); the quadratic logarithm is
\(-t^2/2+\Psi_b(t)/H\). Therefore, for \(|t|\le H^{1/8}\), the exponent
in~\eqref{eq:fourier-full} equals
\begin{equation}
 i\lambda t-\frac{t^2}{2}+\frac{\Psi_b(t)}{H}+E_n(t),
 \label{eq:fourier-exponent}
\end{equation}
where one can bound \(E_n(t)\) by a constant times
\[
 \frac Bn\left(t^2+\frac{t^2(1+|\log|t||)}{H}\right)
 +\frac{n}{B^3}|t|^3.
\]
Here \(t^2|\log|t||\) is assigned its continuous value at zero.
Both prefactors are smaller than every fixed inverse power of \(H\),
including after multiplication by a fixed power of \(H\).

On this interval \(\Psi_b/H=o(1)\) uniformly. Taylor's formula for its
exponential, integrated against the Gaussian, has remainder
\(O(H^{-K-1})\), because
\(\int e^{-t^2/4}|\Psi_b(t)|^{K+1}\,\dd t<\infty\).
The same Gaussian majorant absorbs \(E_n\). By
Lemma~\ref{lem:damping}, the rest of the exact Fourier integral is smaller
than every power of \(1/H\): first there is a Gaussian tail beyond
\(H^{1/8}\), then a bound \(O(B e^{-cB})\), and finally
\(O(B e^{-cn})\). The tails of each approximating Gaussian--logarithmic
integral are equally negligible. Extending those integrals to \(\R\)
gives~\eqref{eq:allorders}. Every estimate used is uniform for bounded
\(\lambda\).
\end{proof}

\subsection{Explicit constants and a finite formula for every order}
\begin{corollary}[Critical coefficient expansion]
\label{cor:critical-coeff}
At \(\beta=\bs\),
\begin{equation}
 u_n(\bs)=\frac{\rho^{-n}}{\sqrt{2\pi}\,nB_n}
 \left[1+\frac{\kappa_1}{H_n}+\frac{\kappa_2}{H_n^2}
 +O(H_n^{-3})\right],
 \label{eq:critical-coeff}
\end{equation}
where
\begin{align}
 \kappa_1&=\frac{2-\gamma-\log2-2b}{4},\label{eq:kappa1}\\
 \kappa_2&=\frac38\left[
 \left(\frac43-\frac{\gamma+\log2}{2}-b\right)^2
 -\frac{\pi^2}{8}-\frac{10}{9}\right].\label{eq:kappa2}
\end{align}
In particular,
\begin{equation}
 u_n(\bs)\sim\frac{\rho^{-n}}{\sqrt{\pi A}\,n^{3/2}\sqrt{\log n}}.
 \label{eq:leading-coefficient}
\end{equation}
\end{corollary}
\begin{proof}
For a standard normal variable \(T\), the first coefficient is
\(\tfrac12\E[T^2(\log|T|-b)]\). Differentiate
\(\E|T|^s=2^{s/2}\Gamma((s+1)/2)/\sqrt\pi\) at \(s=2\) to obtain
\(\E[T^2\log|T|]=1-(\gamma+\log2)/2\).
At the second order, conjugate symmetry replaces the square of
\(\log(-it)-b\) by \((\log|t|-b)^2-\pi^2/4\).
Under the probability measure proportional to \(t^4e^{-t^2/2}\,\dd t\),
\(\log|t|\) has mean
\(4/3-(\gamma+\log2)/2\) and variance \(\pi^2/8-10/9\).
Since \(\E T^4=3\), formula~\eqref{eq:kappa2} follows. Finally,
\(H_n\sim\tfrac12\log n\) gives~\eqref{eq:leading-coefficient}.
\end{proof}

A convenient nonrecursive prescription for every critical coefficient uses
\begin{equation}
 \mathcal M(\alpha)=\frac{1}{2\pi}
 \int_{\R}e^{-t^2/2}(-it)^\alpha\,\dd t
 =\frac{2^{(\alpha-1)/2}}{\pi}
 \Gamma\!\left(\frac{\alpha+1}{2}\right)
 \cos\!\left(\frac{\pi\alpha}{2}\right)
 \quad(\Re\alpha>-1).
 \label{eq:Mellin}
\end{equation}
Differentiation under the integral is justified on compact subsets of that
half-plane. Since \(t^{2k}=(-1)^k(-it)^{2k}\),
\begin{equation}
 G_k^{(b)}(0)=\left.
 \frac{(-1)^k}{2^k k!}(\partial_\alpha-b)^k\mathcal M(\alpha)
 \right|_{\alpha=2k}.
 \label{eq:all-coefficient-formula}
\end{equation}
Thus each coefficient reduces to finitely many gamma derivatives and
ordinary arithmetic. This formula is not an assertion about the growth of
these coefficients as \(k\to\infty\).

\section{A universal cutoff law and the minimal action budget}
\label{sec:cutoff}

\subsection{A quantitative local limit for the truncated action sum}
In this section \(M\) is an integer, \(\ell_n=M/N_n\), and
\(\ell_n\) stays in a compact subset of \((0,\infty)\). Define
\begin{align}
 \Lambda_n&=n\beta_n\sum_{j>M}a_j,\notag\\
 \xi_n&=\frac{n-n\beta_n\sum_{j\le M}ja_j}{B_n},
 &v_n&=\frac{n\beta_n\sum_{j\le M}j^2a_j}{B_n^2}.
 \label{eq:truncated-parameters}
\end{align}
The notation \(v_n\) denotes a normalized variance, not the local inverse
correction \(v(t,r)\) from Section~\ref{sec:chart}.

\begin{lemma}[Truncated Gaussian local approximation]
\label{lem:trunc-local}
Uniformly for bounded \(\lambda\) and the stated \(\ell_n\)-range,
\begin{equation}
 B_n\Prob(S_{n,M}=n)
 =v_n^{-1/2}\ph\!\left(\frac{\xi_n}{\sqrt{v_n}}\right)
 +O(H_n^{-3/2}).
 \label{eq:trunc-local}
\end{equation}
Writing \(\ell=\ell_n\) and
\begin{equation}
 d_n=\nu+\log\ell-\frac12\log H_n,
 \label{eq:d}
\end{equation}
we also have
\begin{align}
 \Lambda_n&=\frac{1}{2\ell^2}+o(H_n^{-K}),\notag\\
 \xi_n&=-\lambda+\frac{1}{\ell\sqrt{H_n}}+O(N_n^{-1}),\notag\\
 v_n&=1+\frac{d_n}{H_n}+o(H_n^{-K})
 \label{eq:parameter-expansions}
\end{align}
for each fixed \(K\), with the little-oh statements uniform on the same
compact sets.
\end{lemma}
\begin{proof}
For \(M\ge J_0\), elementary tail estimates and the harmonic-number
expansion give
\begin{align}
 \sum_{j>M}a_j&=\frac{a}{2M^2}+O(M^{-3}),\notag\\
 \sum_{j>M}ja_j&=\frac aM+O(M^{-2}),\notag\\
 \sum_{j\le M}j^2a_j&=a(\log M+\nu)+O(M^{-1}).
 \label{eq:tail-moments}
\end{align}
Also \(\sum_{j\le M}j^3a_j=O(M)\).
The exact centered Fourier exponent for \(S_{n,M}\), at frequency \(t/B\),
is
\[
 -i\xi_n t-\frac{v_n t^2}{2}+E_{n,M}(t),\qquad
 |E_{n,M}(t)|\le
 \frac{n\beta_n|t|^3}{6B^3}\sum_{j\le M}j^3a_j
 \le C|t|^3H^{-3/2}.
\]
The elementary inequality used here is
\(|e^{ix}-1-ix+x^2/2|\le |x|^3/6\) for real \(x\).
Expand on \(|t|\le H^{1/8}\) and use the same Gaussian damping and
negligible complementary tails as in Theorem~\ref{thm:allorders}.
The integral error is \(O(H^{-3/2})\). Fourier inversion of the quadratic
exponent gives~\eqref{eq:trunc-local}.
Substitution of \(M=\ell N\), \(B=N\sqrt H\), and
\(\log N=H-\tfrac12\log H\) in~\eqref{eq:tail-moments} proves
\eqref{eq:parameter-expansions}. Errors involving \(1/N\) or \(B/n\)
are smaller than every fixed power of \(1/H\).
\end{proof}

\subsection{Coupling--cutoff correction and cancellation of prefix data}
Define a real function by the convergent integral
\begin{equation}
 \cJ(\lambda)=\frac{1}{2\pi}\int_{\R}
 e^{i\lambda t-t^2/2}\frac{t^2}{2}\log(-it)\,\dd t.
 \label{eq:J}
\end{equation}
In particular,
\(\cJ(0)/\ph(0)=(2-\gamma-\log2)/4\).

\begin{theorem}[Universal moving-window cutoff correction]
\label{thm:cutoff-window}
Let
\(R_n(M;\lambda)=u_n^{[M]}(\beta_n)/u_n(\beta_n)\).
For \(\lambda\) in a fixed compact subset of \(\R\) and
\(\ell=\ell_n=M/N_n\) in a fixed compact subset of \((0,\infty)\),
\begin{align}
 R_n(M;\lambda)=e^{-1/(2\ell^2)}
 \Bigg[1&+\frac{\lambda}{\ell\sqrt H}\notag\\
 &+\frac{\lambda^2-1}{2H}
 \left(\gamma-\frac32+\log\ell-\frac12\log H+\frac1{\ell^2}\right)
 -\frac{\cJ(\lambda)}{H\ph(\lambda)}
 +O\!\left(\frac{\log H}{H^{3/2}}\right)\Bigg].
 \label{eq:cutoff-window}
\end{align}
In particular, at critical coupling,
\begin{equation}
 R_n(M;0)=e^{-1/(2\ell^2)}
 \left[1+\frac{\log(2H/\ell^2)+1-\gamma-2/\ell^2}{4H}
 +O(H^{-3/2})\right].
 \label{eq:critical-cutoff}
\end{equation}
Apart from the scale \(A\) entering \(N,H\), these displayed corrections
are unchanged by every finite perturbation allowed in~\eqref{eq:weights}.
\end{theorem}
\begin{proof}
By the exact Poisson identity,
\[
 R_n(M;\lambda)=e^{-\Lambda_n}
 \frac{\Prob(S_{n,M}=n)}{\Prob(S_n=n)}.
\]
The full denominator, multiplied by \(B\), is
\(\ph(\lambda)+G_1^{(b)}(\lambda)/H+O(H^{-2})\).
Its leading term is bounded away from zero on a compact \(\lambda\)-set.
Taylor expansion of the Gaussian term in Lemma~\ref{lem:trunc-local} gives
\begin{align}
 \frac{v_n^{-1/2}\ph(\xi_n/\sqrt{v_n})}{\ph(\lambda)}
 =1+\frac{\lambda}{\ell\sqrt H}
 +\frac{\lambda^2-1}{2H}\left(d_n+\frac1{\ell^2}\right)
 +O\!\left(\frac{\log H}{H^{3/2}}\right).
 \label{eq:gaussian-ratio-expansion}
\end{align}
Indeed, the displacement of the mean is \(1/(\ell\sqrt H)\) and the
variance displacement is \(d_n/H=O(\log H/H)\); retaining the square of
the former produces the \(1/\ell^2\) term. The mixed term has size
\(O(\log H/H^{3/2})\), and the square of the variance displacement is
smaller than that order.

Differentiating the Gaussian Fourier transform twice shows
\begin{equation}
 G_1^{(b)}(\lambda)
 =\cJ(\lambda)+\frac b2(\lambda^2-1)\ph(\lambda).
 \label{eq:b-cancellation}
\end{equation}
Divide~\eqref{eq:gaussian-ratio-expansion} by the normalized full
probability, multiply by \(e^{-\Lambda_n}\), and use
\(\nu-b=\gamma-3/2\). This cancels \(b\) exactly and proves
\eqref{eq:cutoff-window}.

When \(\lambda=0\), odd first-order mean terms and the mixed
mean--variance term vanish. The remaining Taylor error is
\(O((\log H)^2/H^2)=o(H^{-3/2})\), in addition to the local-limit error
\(O(H^{-3/2})\). Finally,
\[
 -\frac{d_n}{2}-\frac1{2\ell^2}-\kappa_1
 =\frac{\log(2H/\ell^2)+1-\gamma-2/\ell^2}{4},
\]
which proves~\eqref{eq:critical-cutoff} with its improved error.
\end{proof}

\begin{remark}[What the universality statement does not say]
Finite-prefix changes alter \(A,b,\rho\), the exact inverse chart, and the
full-sector coefficients \(\kappa_k\). They do not disappear from the
entire problem. The theorem identifies the cancellation in the specific
normalized cutoff ratio through the orders displayed. It makes no
all-orders prefix-independence claim.
\end{remark}

\subsection{Sharp scale, necessity, and an optimally selected budget}
\begin{corollary}[Complete cutoff trichotomy]
\label{cor:trichotomy}
At critical coupling,
\begin{equation}
 R_n(M;0)\longrightarrow
 \begin{cases}
 0,&M/N_n\longrightarrow0,\\
 e^{-1/(2\ell^2)},&M/N_n\longrightarrow\ell\in(0,\infty),\\
 1,&M/N_n\longrightarrow\infty.
 \end{cases}
 \label{eq:trichotomy}
\end{equation}
Moreover, \(R_n(M;0)\to1\) if and only if \(M/N_n\to\infty\).
The same conclusions hold in the moving window for any bounded sequence
\(\lambda=\lambda_n\).
\end{corollary}
\begin{proof}
For finite positive \(\ell\), apply Theorem~\ref{thm:cutoff-window}.
For \(M/N\to0\), monotonicity in \(M\) bounds the limsup by
\(e^{-1/(2\varepsilon^2)}\) for each fixed \(\varepsilon>0\); let
\(\varepsilon\downarrow0\). For \(M/N\to\infty\), the liminf is at least
\(e^{-1/(2L^2)}\) for each fixed \(L>0\); let \(L\to\infty\) and use
\(R_n\le1\). If \(M/N\) fails to tend to infinity, a subsequence is bounded
above by a fixed \(L\), and its ratio has limsup at most
\(e^{-1/(2L^2)}<1\). Uniformity in bounded \(\lambda\) gives the last
assertion. Integer parts in the comparison cutoffs have no effect.
\end{proof}

\Needspace{10\baselineskip}
\begin{theorem}[Two-term minimal action budget]
\label{thm:min-budget}
Fix \(0<\eta<1\), let \(d_\eta=-\log(1-\eta)\), and define
\[
 M_{\min}(n,\eta)=\min\{M\ge1:R_n(M;0)\ge1-\eta\}.
\]
Then
\begin{equation}
 M_{\min}(n,\eta)=\frac{N_n}{\sqrt{2d_\eta}}
 \left[1-
 \frac{\log(4d_\eta H_n)+1-\gamma-4d_\eta}{8d_\eta H_n}
 +O(H_n^{-3/2})\right].
 \label{eq:min-budget}
\end{equation}
The error statement is for fixed \(\eta\); it is also uniform when
\(\eta\) stays in a compact subset of \((0,1)\).
\end{theorem}
\begin{proof}
Existence follows from exact equality of the coefficients for \(M\ge n\).
The leading profile and monotonicity first imply
\(M_{\min}/N\to\ell_0=(2d_\eta)^{-1/2}\).
Take logarithms in~\eqref{eq:critical-cutoff}. The square of the correction
is \(o(H^{-3/2})\), so
\[
 \log R_n(M;0)=-\frac1{2\ell^2}+\frac{C(\ell,H)}{H}+O(H^{-3/2}),
 \quad C(\ell,H)=\frac{\log(2H/\ell^2)+1-\gamma-2/\ell^2}{4}.
\]
Writing \(\ell=\ell_0(1+z/H)\), the first term is
\(-d_\eta+2d_\eta z/H+O(z^2/H^2)\). The threshold therefore occurs at
\[
 z=-\frac{\log(4d_\eta H)+1-\gamma-4d_\eta}{8d_\eta},
\]
up to an error \(O(H^{-1/2})\). More explicitly, perturbing the proposed
\(\ell\) by \(\pm C_0H^{-3/2}\), with \(C_0\) sufficiently large, gives
opposite strict inequalities for the threshold because the derivative
of \(-1/(2\ell^2)\) is bounded below near \(\ell_0\). Monotonicity traps
the first crossing. Rounding to the next integer changes \(\ell\) by
\(O(N^{-1})=o(H^{-K})\) for every fixed \(K\).
\end{proof}

The theorem supplies an asymptotic optimizer, not an unconditional finite
error guarantee obtained by discarding its big-oh term. The next section
provides such a guarantee separately.

\section{A finite error certificate and proof-producing computation}
\label{sec:certificate}

\subsection{An elementary anti-concentration bound}
The following estimate needs less than the exact-tail assumption. Suppose
only that \(a_j\ge0\), \(a_1>0\), \(\sum a_j<\infty\), and \(\beta>0\).
The finite sums below are always defined. Let \(S_n\) be the compound Poisson
sum of Proposition~\ref{prop:poisson}.

\begin{proposition}[Explicit bound on every point probability]
\label{prop:anti-concentration}
For any integer \(J\ge1\) and \(0<\vartheta\le1\), put
\begin{align}
 V_J&=\sum_{j\le J}j^2a_j,\notag\\
 I_{n,J,\vartheta}
 &=\frac{1}{\sqrt{2\pi n\beta(1-\vartheta^2/12)V_J}}
 +\exp\!\left(-\frac{2n\beta a_1\vartheta^2}{\pi^2J^2}\right).
 \label{eq:I-bound}
\end{align}
Then \(\sup_{k\in\mathbb Z}\Prob(S_n=k)\le I_{n,J,\vartheta}\).
\end{proposition}
\begin{proof}
Taking absolute values in Fourier inversion gives
\[
 \Prob(S_n=k)\le\frac1\pi\int_0^\pi
 \exp\left\{-n\beta\sum_{j\ge1}a_j(1-\cos j\theta)\right\}\,\dd\theta.
\]
If \(0\le\theta\le\vartheta/J\), then
\[
 1-\cos(j\theta)\ge\frac{j^2\theta^2}{2}
 \left(1-\frac{\vartheta^2}{12}\right)\quad(j\le J).
\]
Extend the resulting Gaussian integral to infinity to obtain the first
term of~\eqref{eq:I-bound}. On the rest of \([0,\pi]\), the single weight
\(a_1\) gives
\(\sum a_j(1-\cos j\theta)\ge a_1(1-\cos\theta)
\ge2a_1\theta^2/\pi^2\).
Bounding the integrand by its value at \(\vartheta/J\) and the interval
length by \(\pi\) gives the second term. The estimates are independent of
\(k\).
\end{proof}

\subsection{A positive coefficient recurrence with a certified tail inequality}
For fixed \(n,M,\beta\), define
\begin{equation}
 P_{n,M}=e^{-n\beta F(1)}[z^n]e^{n\beta F_M(z)},\qquad
 \Lambda=n\beta\sum_{j>M}a_j.
 \label{eq:P-certificate}
\end{equation}
This normalization uses the full \(F(1)\), not \(F_M(1)\). Accordingly,
\(P_{n,M}\) is the probability of \(S_n=n\) and no omitted action occurring.

\begin{theorem}[Finite coefficient-retention certificate]
\label{thm:finite-certificate}
For every \(n,M,J\ge1\), \(\beta>0\), and \(0<\vartheta\le1\),
\begin{equation}
 \boxed{\quad
 \frac{P_{n,M}}{P_{n,M}+\Lambda I_{n,J,\vartheta}}
 \le\frac{u_n^{[M]}(\beta)}{u_n(\beta)}\le1.
 \quad}
 \label{eq:finite-certificate}
\end{equation}
All quantities on the left are explicit. In the exact-tail model at
critical coupling and \(M\ge J_0\), one may replace \(\Lambda\) by the
upper bound \(nA/(2M^2)\).
\end{theorem}
\begin{proof}
Let \(Z_{>M}=\sum_{j>M}Z_j\). The missed probability satisfies
\begin{align*}
 \Prob(S_n=n,Z_{>M}\ge1)
 &\le\E[Z_{>M}\mathbf1_{\{S_n=n\}}]\\
 &=\sum_{j>M}n\beta a_j\Prob(S_n=n-j)\\
 &\le\Lambda I_{n,J,\vartheta}.
\end{align*}
The middle equality is the elementary Poisson shift identity
\(\E[Z_jh(S_n)]=n\beta a_j\E[h(S_n+j)]\), proved by shifting the index
in the Poisson mass function. Nonnegative terms justify summing over the
countable family. In fact, terms with \(j>n\) vanish, though using the
whole tail is a convenient upper bound. Since
\[
 \Prob(S_n=n)=P_{n,M}+\Prob(S_n=n,Z_{>M}\ge1),
\]
and the desired coefficient ratio is \(P_{n,M}/\Prob(S_n=n)\), the result
follows. Finally,
\(\sum_{j>M}j^{-3}\le\int_M^\infty x^{-3}\,\dd x=1/(2M^2)\).
\end{proof}

The coefficient \(P_{n,M}\) is computable by a recurrence containing no
subtractions. With \(p_0=e^{-n\beta F(1)}\), define
\begin{equation}
 p_k=\frac{n\beta}{k}
 \sum_{j=1}^{\min(k,M)}j a_j p_{k-j}\quad(1\le k\le n).
 \label{eq:positive-recurrence}
\end{equation}
Then \(p_n=P_{n,M}\). This follows by differentiating
\(\sum p_kz^k=\exp(n\beta(F_M(z)-F(1)))\).
The arithmetic-operation count is \(O(nM)\), and an implementation needs
only the preceding \(M\) values, hence \(O(M)\) working storage.
These are arithmetic-operation bounds, not bit-complexity or rounding-error
claims. The reference script keeps all coefficients for simplicity.

For a literal machine certificate, evaluate a lower enclosure \(P_-\le P\)
and upper enclosures \(\Lambda_+\ge\Lambda\), \(I_+\ge I\), all nonnegative.
Then
\[
 \frac{P_-}{P_-+\Lambda_+I_+}
\]
is a valid lower retention bound, with its final division also rounded
downward. The mpmath diagnostics supplied here do not perform that directed
rounding. Thus Theorem~\ref{thm:finite-certificate} is a proved mathematical
certificate, whereas the recorded floating-point values illustrate it rather
than instantiate a proof-assistant or interval-arithmetic certificate.

\subsection{The finite certificate attains the sharp order of the action budget}
The Fourier majorant can be made asymptotically sharp enough to preserve the
\(\sqrt n\) scale. Let
\begin{equation}
 J_n=\max\left(1,\left\lfloor\frac{\sqrt n}{(\log n)^{3/2}}\right\rfloor\right),
 \qquad \vartheta_n=(\log n)^{-1/2},
 \label{eq:J-choice}
\end{equation}
for \(n\) sufficiently large that \(\vartheta_n\le1\).

\begin{corollary}[Asymptotic cost of the finite certificate]
\label{cor:certificate-budget}
At critical coupling,
\[
 I_{n,J_n,\vartheta_n}\sim\frac{1}{\sqrt{2\pi}B_n}.
\]
If \(M/N_n\to\ell\in(0,\infty)\), and \(t=1/(2\ell^2)\), the lower
bound in~\eqref{eq:finite-certificate} tends to
\begin{equation}
 \frac{1}{1+t e^t}.
 \label{eq:certificate-profile}
\end{equation}
Consequently, for a fixed tolerance \(0<\eta<1\), every fixed constant
\begin{equation}
 \ell>\frac{1}{\sqrt{2W_0(\eta/(1-\eta))}}
 \label{eq:certificate-budget}
\end{equation}
is sufficient for the finite inequality eventually to certify
\(R_n(\lceil\ell N_n\rceil;0)\ge1-\eta\).
The ratio of this threshold constant to the optimal asymptotic constant is
\begin{equation}
 \sqrt{\frac{-\log(1-\eta)}{W_0(\eta/(1-\eta))}}
 \longrightarrow1\qquad(\eta\downarrow0).
 \label{eq:certificate-efficiency}
\end{equation}
The limits here are first \(n\to\infty\) for fixed \(\eta\), then
\(\eta\downarrow0\).
\end{corollary}
\begin{proof}
We have \(V_{J_n}\sim a\log J_n\sim aH_n\), so the Gaussian term
in~\eqref{eq:I-bound} is asymptotic to \(1/(\sqrt{2\pi}B_n)\).
The second term is \(\exp(-c(\log n)^2(1+o(1)))=o(B_n^{-1})\).
Theorem~\ref{thm:cutoff-window} and the local limit give
\(P_{n,M}\sim e^{-t}/(\sqrt{2\pi}B_n)\), while \(\Lambda\to t\).
Substitution proves~\eqref{eq:certificate-profile}.
Its value exceeds \(1-\eta\) precisely when
\(te^t<\eta/(1-\eta)\), proving~\eqref{eq:certificate-budget}.
Finally, both \(-\log(1-\eta)\) and \(W_0(\eta/(1-\eta))\) are
asymptotic to \(\eta\).
\end{proof}

This comparison does not make the elementary majorant optimal at every
fixed tolerance. It proves that an actual finite inequality, rather than an
uncontrolled asymptotic substitution, operates on the correct action scale.
For a prescribed finite \(n\), one evaluates the inequality itself.

\section{How the finite-action folds approach the marginal boundary}
\label{sec:fold}

\subsection{A second convergent logarithmic chart}
Fix \(\beta=\bs\). The finite map
\(q=z e^{-\bs F_M(z)}\) has a unique positive critical point
\(z_M=e^{v_M}>1\), characterized by
\begin{equation}
 \bs\sum_{j\le M}j a_j e^{jv_M}=1.
 \label{eq:finite-fold-equation}
\end{equation}
Indeed, the left side is strictly increasing in \(v_M\), is below one at
zero, and tends to infinity. Put
\[
 L_M=\log M+\nu,\qquad \rho_M=e^{v_M-\bs F_M(e^{v_M})}.
\]
The reuse of the letter \(L\) with subscript \(M\) should not be confused
with the inverse variable \(L\) in~\eqref{eq:log-generators}.
Define entire functions
\begin{align}
 K(s)&=\sum_{k\ge2}\frac{s^k}{k!(k-1)}
 =\int_0^1\frac{e^{sx}-1-sx}{x^2}\,\dd x,\notag\\
 Q(s)&=\sum_{k\ge3}\frac{s^k}{k!(k-2)}
 =\int_0^1\frac{e^{sx}-1-sx-(sx)^2/2}{x^3}\,\dd x.
 \label{eq:KQ}
\end{align}
The integrands have removable values at zero.

\begin{theorem}[Logarithmic fold chart, drift, and curvature]
\label{thm:finite-fold}
There is a holomorphic \(S(r)\) near zero, uniquely determined by
\begin{equation}
 S+rK(S)=r,\qquad S(0)=0.
 \label{eq:S-equation}
\end{equation}
It has a convergent Taylor series, beginning
\begin{equation}
 S(r)=r-\frac12r^3-\frac1{12}r^4+O(r^5).
 \label{eq:S-expansion}
\end{equation}
For the finite-action critical point,
\begin{equation}
 v_M=\frac1M S(1/L_M)+O\!\left(\frac{1}{M^2L_M}\right).
 \label{eq:finite-fold-v}
\end{equation}
Its critical value satisfies
\begin{align}
 \log\frac{\rho_M}{\rho}
 =\frac A{M^2}\left[
 \frac12+S-\frac{L_M}{2}S^2-Q(S)\right]+O(M^{-3}),
 \qquad S=S(1/L_M),
 \label{eq:finite-fold-rho}
\end{align}
and hence
\begin{equation}
 \log\frac{\rho_M}{\rho}
 =\frac A{2M^2}\left[1+\frac1{L_M}+O(L_M^{-3})\right]+O(M^{-3}).
 \label{eq:finite-fold-rho-simple}
\end{equation}
Finally, the positive curvature is
\begin{equation}
 \kappa_M:=\bs\sum_{j\le M}j^2a_j e^{jv_M}
 =A\left[L_M+K'\bigl(S(1/L_M)\bigr)\right]+O(M^{-1}).
 \label{eq:finite-fold-curvature}
\end{equation}
\end{theorem}
\begin{proof}
The implicit function theorem applied to \(S+rK(S)-r\) at zero gives
\(S(r)\), since the derivative with respect to \(S\) is one. Coefficient
comparison gives~\eqref{eq:S-expansion}.

Write \(s=Mv_M\). First observe that \(s=O(1/L_M)\): by positivity,
\(e^{jv}-1\ge jv\), so the critical equation implies
\[
 v_M\sum_{j\le M}j^2a_j\le\sum_{j>M}ja_j.
\]
The moment estimates in~\eqref{eq:tail-moments} give the assertion.
Expand the left side of~\eqref{eq:finite-fold-equation} after subtracting
its value at zero. The linear term is \(aL_Ms/M+O(M^{-2})\).
For the exact tail, the nonlinear part, multiplied by \(M/a\), is
\[
 \frac1M\sum_{j\le M}
 \frac{e^{s(j/M)}-1-s(j/M)}{(j/M)^2}=K(s)+O(M^{-1}),
\]
uniformly for \(s\) in a bounded complex neighborhood of zero. This is a
Riemann sum for a function analytic jointly in \(s,x\) after removing
its value at \(x=0\); its first \(x\)-derivative is uniformly bounded
on the relevant compact set. The finitely many perturbations contribute
\(O(M^{-1})\) after the same scaling. The omitted first moment is
\(a/M+O(M^{-2})\). We obtain
\begin{equation}
 L_Ms+K(s)=1+O(M^{-1}).
 \label{eq:fold-approx-equation}
\end{equation}
The derivative on the left is \(L_M+O(1)\). Comparison with the exact
solution \(S(1/L_M)\) gives \(s-S=O(1/(ML_M))\), proving
\eqref{eq:finite-fold-v}.

For the critical-value drift, Taylor expansion at \(z=1\) in the
logarithmic coordinate gives
\begin{align*}
 \log\frac{\rho_M}{\rho}
 ={}&\bs(F(1)-F_M(1))
 +v_M\left(1-\bs\sum_{j\le M}ja_j\right)\\
 &-\frac{\bs v_M^2}{2}\sum_{j\le M}j^2a_j
 -\bs\sum_{j\le M}a_j
 \left(e^{jv_M}-1-jv_M-\frac{(jv_M)^2}{2}\right).
\end{align*}
The four terms, multiplied by \(M^2/A\), are respectively
\(1/2+O(M^{-1})\), \(s+O(M^{-1})\),
\(-L_Ms^2/2+O(M^{-1})\), and \(-Q(s)+O(M^{-1})\).
The last estimate follows from the analogous analytic Riemann sum with
an \(x^{-3}\) denominator. Replacing \(s\) by \(S\) changes this scaled
expression by at most \(O(M^{-1})\), proving~\eqref{eq:finite-fold-rho}.
To expand it, use \(L_MS=1-K(S)\); the bracket equals
\[
 \frac12+\frac S2+\frac{SK(S)}2-Q(S)
 =\frac12+\frac{1}{2L_M}+O(L_M^{-3}).
\]
This also explains the absence of a \(L_M^{-2}\) term.

Finally,
\[
 \sum_{j\le M}j^2a_j(e^{jv_M}-1)
 =a\int_0^1\frac{e^{sx}-1}{x}\,\dd x+O(M^{-1})
 =aK'(s)+O(M^{-1}),
\]
by a third analytic Riemann-sum estimate; finite-prefix changes are
\(O(v_M)\). Add the unperturbed truncated second moment and replace
\(s\) by \(S\). This proves~\eqref{eq:finite-fold-curvature}.
\end{proof}

\subsection{The analytic fold law and why the two limits differ}
For each fixed \(M\), put \(U_M=\bs F_M(e^{v_M})\). On the positive
small-solution sheet, as \(q\uparrow\rho_M\),
\begin{equation}
 U_{\bs}^{[M]}(q)=U_M-
 \sqrt{\frac{2\log(\rho_M/q)}{\kappa_M}}
 +O\bigl(\log(\rho_M/q)\bigr).
 \label{eq:local-finite-fold}
\end{equation}
To see this directly, write \(v=\log z\). The derivative of
\(\log q=v-\bs F_M(e^v)\) vanishes at \(v_M\), and its second derivative
there is \(-\kappa_M\). The derivative of \(\bs F_M(e^v)\) at the same
point is one. The analytic square-root inversion of this nondegenerate
quadratic zero gives~\eqref{eq:local-finite-fold}.

Every finite model therefore has an ordinary analytic fold. Nevertheless,
\(\kappa_M\sim A\log M\) diverges, \(v_M\sim1/(M\log M)\), and the
limiting boundary has the logarithmic inverse of Section~\ref{sec:chart}.
At a cutoff \(M\sim\ell N_n\), the critical-value drift satisfies
\begin{equation}
 n\log(\rho_M/\rho)\longrightarrow\frac1{2\ell^2}.
 \label{eq:fold-profile-link}
\end{equation}
Its exponential is exactly the leading suppression factor in the cutoff
law. This is an explanatory compatibility check, not a substitute for a
uniform coefficient theorem: fixed-\(M\) singularity analysis alone would
not justify letting \(M\) grow with \(n\). Theorem~\ref{thm:cutoff-window}
was proved independently by uniform Fourier estimates.

\section{Reproducible verification and numerical examples}
\label{sec:verification}

\subsection{What was checked exactly}
The package contains two executable Python programs and their recorded JSON
outputs. The main program performs 102 exact checks: 96 comparisons of a
rational exponential-series recurrence with independent partition
enumeration; four coefficients of the reduced Lambert-core equation; the
cubic cancellation caused by increasing \(a_1\) by \(1/2\); and the algebraic
cancellation of \(b\) in the critical cutoff correction. The partition tests
use \(n=1,\ldots,12\), \(\beta=1/3,2/3\), two prefixes, and both a fixed
small cutoff and the full degree cutoff. They check the coefficient before
the common factor \(1/n\).

The supplementary program performs nine further exact checks: eight
coefficients of the fold equation~\eqref{eq:S-equation}, through degree
seven, and the finite-prefix cancellation for general \(\lambda\).
Thus 111 finite exact checks passed in the recorded runs. The independent
fold computation gives
\[
 S(r)=r-\frac{r^3}{2}-\frac{r^4}{12}
 +\frac{35r^5}{72}+\frac{33r^6}{160}
 -\frac{1013r^7}{1800}+O(r^8).
\]
These checks do not verify infinitely many coefficients or replace any
analytic proof above.

\subsection{Inverse accuracy and sector coefficients}
For the pure weights \(a_j=j^{-3}\),
\[
 A=0.60792710185402662866\ldots,\qquad b=\frac32,
\]
and the first two normalized sector constants are
\[
 \kappa_1=-0.56759071136536954251\ldots,\qquad
 \kappa_2=-0.63819423854137909958\ldots.
\]
Direct quadrature of the Gaussian--logarithmic moments agrees with the
closed formulas within \(2\times10^{-50}\) in the fifty-digit run.
A separate comparison at \(n=256\) checks the Fourier calculation against
the positive coefficient recurrence, for both the full coefficient and
\(M=\lfloor N_n\rfloor\). The relative discrepancies are below
\(10^{-48}\). These are floating-point agreement figures, not rigorous
error enclosures.

\begin{table}[htbp]
\centering\small
\begin{tabular}{@{}rrrr@{}}
\toprule
\(\epsilon\)&\(\tau\), numerical root&Core relative error&Three-term relative error\\
\midrule
\(10^{-4}\)&\(7.146046748\times10^{-3}\)&\(1.0020\times10^{-4}\)&\(2.7144\times10^{-11}\)\\
\(10^{-8}\)&\(5.388914675\times10^{-5}\)&\(4.1471\times10^{-7}\)&\(2.9478\times10^{-18}\)\\
\(10^{-16}\)&\(3.972838893\times10^{-9}\)&\(1.6274\times10^{-11}\)&\(3.0045\times10^{-31}\)\\
\bottomrule
\end{tabular}
\caption{Pure-tail inverse diagnostics. The core is \(t\); the three-term
approximation retains \(t,t^2,t^3\) in~\eqref{eq:pure-inverse}.}
\label{tab:inverse}
\end{table}

\begin{table}[htbp]
\centering\small
\begin{tabular}{@{}rrrr@{}}
\toprule
\(n\)&\(H_n\)&\(\sqrt{2\pi}\,nB_n\rho^n u_n\)&\(1+\kappa_1/H_n+\kappa_2/H_n^2\)\\
\midrule
\(10^4\)&5.1785861&0.87607723&0.86659913\\
\(10^8\)&10.1186819&0.93893996&0.93767354\\
\(10^{16}\)&19.6611530&0.96965125&0.96948041\\
\bottomrule
\end{tabular}
\caption{Pure-tail sector coefficients and their first two corrections.
Convergence in the logarithmic variable is slow even at very large \(n\).}
\label{tab:coefficients}
\end{table}

The inverse chart is genuinely convergent; the outer coefficient expansion
is only asymptotic. The sharp contrast between Tables~\ref{tab:inverse}
and~\ref{tab:coefficients} is therefore expected. A large value of \(n\)
need not make \(1/H_n\) extremely small.

\subsection{Cutoff, prefix, and moving-window diagnostics}
The cutoff calculations use the exact integer \(M\), not its unrounded
nominal value, in \(\ell=M/N_n\). Table~\ref{tab:cutoff} uses
\(M=\lfloor N_n\rfloor\). The complete JSON file additionally contains
nominal \(\ell=1/2,2\), both prefixes, and more digits.

\begin{table}[htbp]
\centering\small
\begin{tabular}{@{}lrrrr@{}}
\toprule
Prefix&\(n\)&Numerical ratio&Leading profile&First correction\\
\midrule
\(\delta_1=0\)&\(10^4\)&0.62498322&0.59889287&0.62013748\\
\(\delta_1=0\)&\(10^8\)&0.62820473&0.60645535&0.62788283\\
\(\delta_1=0\)&\(10^{16}\)&0.62290548&0.60653066&0.62268465\\
\(\delta_1=1/2\)&\(10^4\)&0.63186838&0.60403388&0.62574393\\
\(\delta_1=1/2\)&\(10^8\)&0.62794894&0.60644250&0.62795760\\
\(\delta_1=1/2\)&\(10^{16}\)&0.62274259&0.60653065&0.62274334\\
\bottomrule
\end{tabular}
\caption{Retention ratios at critical coupling. Each model uses its own
\(A,H_n,N_n\). The first correction is~\eqref{eq:critical-cutoff}.}
\label{tab:cutoff}
\end{table}

At \(n=10^8\), pure weights, and \(M=7796\), the moving-window ratios are
approximately \(0.45615931\) for \(\lambda=-1\) and \(0.75444199\) for
\(\lambda=1\). Formula~\eqref{eq:cutoff-window} gives respectively
\(0.49001174\) and \(0.74751311\), versus a common leading value
\(0.60645535\). These values illustrate the sizable
\(H^{-1/2}\) coupling correction. They also show why an asymptotic correction
must not be advertised as a finite certified bound: at this order its
accuracy need not be symmetric in \(\lambda\), and its remainder constant
has not been optimized.

For the mathematical finite certificate, the main program evaluates
\eqref{eq:finite-certificate} with \(\vartheta=1\),
\(J=\lfloor\sqrt n/\log(n+2)\rfloor\), and \(M=\lceil3N_n\rceil\).
For \(n=64,M=19\), it gives the numerical lower value \(0.89482402\),
while the numerically computed ratio is \(0.99453947\). For
\(n=256,M=38\), the corresponding values are \(0.90564333\) and
\(0.99161984\). Both calculations respect the proved inequality; neither
uses outward rounding.

\begin{table}[htbp]
\centering\small
\begin{tabular}{@{}rrrrr@{}}
\toprule
\(M\)&\(ML_Mv_M\)&\(2M^2\log(\rho_M/\rho)/A\)&Reduced prediction&\(\kappa_M/(AL_M)\)\\
\midrule
100&0.97577094&1.17856522&1.19057584&1.03909771\\
1000&0.99048085&1.13166298&1.13280358&1.01834924\\
10000&0.99469179&1.10170261&1.10181329&1.01065843\\
\bottomrule
\end{tabular}
\caption{Pure-tail finite-fold diagnostics. The reduced prediction uses the
convergent solution \(S(1/L_M)\) in~\eqref{eq:finite-fold-rho}. Its error
is consistent with the separate algebraic remainder in \(1/M\).}
\label{tab:fold}
\end{table}

\subsection{Reproduction and numerical limitations}
The recorded environment used Python 3.13.5, SymPy 1.14.0, and mpmath 1.3.0.
The primary run uses fifty decimal working digits, and the supplementary
run uses forty. The models implemented numerically are
\(a_j=j^{-3}+\delta_1\mathbf1_{\{j=1\}}\), with
\(\delta_1=0,1/2\); the proofs cover the larger class in~\eqref{eq:weights}.
Fourier diagnostics truncate the scaled integration interval and the
analytic series in the integrand. Their tails are not enclosed by interval
arithmetic in the code. Exact finite comparisons and independent recurrence
checks help detect mistakes, but working precision is not a certified number
of correct digits.

From the package root, run
\begin{verbatim}
python -m pip install -r requirements.txt
python code/verify.py
python code/supplementary.py
sh build.sh
\end{verbatim}
For a smaller numerical run, use \texttt{python code/verify.py --quick}.
The article is self-contained and its typesetting does not require rerunning
any computation or accessing the repository. All displayed numerical tables
are recorded outputs rounded for readability.

\section{Interpretation, proof boundaries, and formalization targets}
\label{sec:interpretation}

\subsection{Three uses of asymptotics, with different logical content}
There are three distinct objects. Near \(q=0\), the solution is an ordinary
convergent series, or equivalently an exponential transseries after
\(q=e^{-X}\). Near the critical boundary \(q=\rho\), the inverse is a
convergent analytic function of a finite logarithmic generating grid. Finally,
as the sector index \(n\to\infty\), the coefficients have a Gaussian--logarithmic
asymptotic expansion. Convergence of the second object does not prove
convergence or Borel summability of the third.

Similarly, \(M\to\infty\) and \(n\to\infty\) are not interchangeable
without a scale hypothesis. The finite-action folds exist for each \(M\),
but the exactly marginal boundary has no finite Puiseux chart. Their reduced
logarithmic chart \(S(1/L_M)\) is convergent; the theorem also keeps a
separate, explicitly stated algebraic \(1/M\) remainder. It does not claim
that every term of a joint infinite \(1/M\)-expansion converges.

\subsection{Why the endpoint cutoff is especially efficient}
The exact ratio can be written
\[
 R_n(M;0)=\Prob(\text{no action exceeds }M\mid S_n=n).
\]
At \(M\asymp N_n\), forbidding large actions has a nontrivial Poisson cost
\(e^{-1/(2\ell^2)}\). Its effect on the mean of the normalized sum is only
\(O(H_n^{-1/2})\), because the fluctuation scale is larger by \(\sqrt H\).
Its effect on the normalized variance is \(O(\log H/H)\).
The normal local limit therefore does not add a new leading cost at the
critical center. This yields a pure extreme-value profile, unlike a regime
where individual large actions and the sum fluctuate on the same scale.

This explanation is made quantitative by the mean and variance shifts in
Lemma~\ref{lem:trunc-local}. The cancellation of the prefix in the ratio
requires comparing those shifts with the first correction to the full
probability. It would be missed by retaining only the leading Gaussian.
Thus the universality calculation is more than the leading extreme-value
observation.

\subsection{Dependency audit}
The main proofs have a short, inspectable dependency structure. The exact
coefficient formula uses formal Lagrange inversion and the analytic implicit
function theorem near zero. The critical chart uses the local polylogarithm
identity, another implicit-function argument, and an explicit contraction
bound. The sector and cutoff theorems use Fourier inversion, elementary
damping from positivity, Gaussian integrals, and Taylor estimates. The
finite certificate uses only Fourier absolute values, a cosine inequality,
and the Poisson shift identity. The fold drift uses analytic Riemann sums
and a one-variable implicit equation.

No result depends on an unproved statement in the repository, an assumed
resurgence theorem, a numerical fit, or a global complex singularity
classification. Conversely, the article does not establish such a
classification, an arbitrary-support Hahn-field closure theorem, a signed-weight
analogue, or general transseries realization. This separation is necessary:
a local analytic chart and a formal support statement are not automatically
an asymptotic realization theorem for an unrestricted formal field.

\subsection{A realistic Lean development}
No new Lean file is supplied or claimed to compile. The inspected repository
already has supporting asymptotic-scale, flatness, and well-based-support
work, but that does not formalize the probability and analytic estimates of
this article \cite{repo-index}. A useful implementation order is the following.

First formalize the finite coefficient recurrence, its equality with the
exponential coefficient, monotonicity in the cutoff, and exact equality
through degree \(M\). These are finite algebraic statements and are the
closest match to the exact checks supplied here. Next isolate the Poisson
shift identity, finite-cutoff conditioning identity, and elementary
anti-concentration majorant. Together they prove the finite certificate
without requiring the full asymptotic theory.

The analytic stage should expose the hypotheses of the local chart as
explicit norm bounds and use a quantitative contraction interface. A
separate Fourier stage should package uniform damping, integrable Taylor
remainders, and the large-scale limits. Only after these pieces are formalized
should the prefix-cancellation and minimal-budget theorems be assembled.
This division prevents a verified finite coefficient identity from being
misreported as a verified infinite asymptotic theorem.

\section{Further research questions and concrete targets}
\label{sec:questions}

\paragraph{1. A uniform transition as the stable exponent approaches two.}
For \(a_j=a j^{-1-\alpha}\), let \(\alpha\uparrow2\) while
\((2-\alpha)\log n\) remains bounded. Construct a joint inverse chart and
coefficient expansion that displays the cancellation between the singular
gamma factor and the quadratic analytic term. Determine how the stable
cutoff density crosses over to \(e^{-1/(2\ell^2)}\), and whether the
prefix-independent first correction has a uniform interpolating form.
Separate limits in \(\alpha\) and \(n\) do not answer this question.

\paragraph{2. Slowly varying perturbations of the marginal tail.}
Replace the exact tail by \(a_j\sim a j^{-3}L(j)\), with a specified slowly
varying function \(L\). The sum scale is governed by a truncated second
moment, while the cutoff scale is governed by \(n\sum_{j>M}a_j\).
Find conditions under which their ratio diverges, remains bounded, or has
an explicit logarithmic law. Identify classes of \(L\) that preserve a
finite convergent generating grid, and classes that require genuinely new
inverse notation or de Bruijn conjugates.

\paragraph{3. Higher cutoff corrections and the return of prefix data.}
Compute the full expansion in \(H^{-1/2}\) and \(\log H/H\), with a
systematic Gaussian--logarithmic coefficient calculus. Determine the first
order at which finite-prefix information survives in the normalized ratio,
if it does. The present cancellation establishes only the displayed first
correction; it is not evidence that every higher order is universal.
A useful deliverable would be a symbolic algorithm with an explicit
uniform remainder for each fixed requested order.

\paragraph{4. Large-order growth of the logarithmic sector expansion.}
Analyze \(G_k^{(b)}(\lambda)\) as \(k\to\infty\), uniformly on useful
\(\lambda\)-sets. Determine least-term truncation, exponentially improved
remainders, and possible Borel or accelerated summation. The finite formula
\eqref{eq:all-coefficient-formula} provides a concrete starting point, but
its existence alone does not establish a Gevrey class or resurgence.
The relation between the convergent critical chart and this potentially
divergent outer expansion deserves a separate theorem.

\paragraph{5. Growing coupling windows and moderate deviations.}
Find the largest range \(|\lambda|\le L_n\to\infty\) on which a relative,
rather than additive, version of Theorem~\ref{thm:allorders} holds. Match it
to subcritical one-large-action behavior and supercritical saddle behavior.
The division by \(\ph(\lambda)\) in the cutoff correction makes clear why
compact-window estimates cannot simply be extended without new analysis.

\paragraph{6. Joint precision and action-budget limits.}
Let \(\eta=\eta_n\downarrow0\) with \(n\). Determine the uniform range of
\eqref{eq:min-budget}, and the sharp relation between the asymptotically
minimal budget and the certificate budget. The successive-limit statement
in Corollary~\ref{cor:certificate-budget} does not settle this joint limit.
Develop an adaptive choice of \(J,\vartheta\) that improves the finite
majorant without losing a simple, checkable expression.

\paragraph{7. Interval-certified coefficient extraction and bit complexity.}
Implement the positive recurrence with directed rounding and explicit
special-function enclosures, producing a machine-readable retention
certificate. Then investigate faster coefficient algorithms, block
convolution, and reliable rescaling for very small normalized coefficients.
Compare their bit complexity with the arithmetic bound \(O(nM)\). A useful
formal target is a small checker for the final inequality, independent of
the heuristic algorithm used to choose \(M\).

\paragraph{8. Vector feedback and noncommensurate actions.}
Replace the scalar feedback by a positive multitype system or a finite
number of coupled marginal tails. Determine whether the covariance matrix
has a logarithmically divergent subspace and how the extreme-action budget
interacts with the remaining finite-variance directions. Noncommensurate
actions require a replacement for integer coefficient extraction; the
appropriate object may be a Laplace measure rather than a single power
series. This is not a routine change of notation.

\paragraph{9. Signed weights, complex saddles, and analytic continuation.}
Allow finite or controlled infinite sign changes while preserving a chosen
analytic small solution. Determine which cutoff conclusions survive without
coefficient monotonicity or a positive Poisson representation. Even for the
positive model, describe global continuation of the Lambert-core chart and
the singular geometry encountered away from the real marginal boundary.
The local pole \(w=1\) in coefficient functions need not by itself identify
an actual singularity of the full solution in a different parameter regime.

\paragraph{10. Formalization of the finite-to-asymptotic bridge.}
Prove the finite certificate first in Lean, then connect it to a separately
formalized marginal local limit theorem. An important interface question is
how to represent compact-uniform Fourier bounds and analytic logarithm
branches so that the analytic hypotheses remain visible in downstream
statements. The final artifact should distinguish a certified numerical
retention ratio from a formally proved asymptotic budget law.

\section{Conclusion}
The marginal \(j^{-3}\) feedback class lies between ordinary finite-action
folds and nonintegral stable criticality. Its logarithm is not a defect to
be suppressed: an exact Lambert--\(W_{-1}\) core turns it into a convergent,
explicitly controlled inverse chart. At the coefficient level, a different
logarithmic expansion describes a Gaussian sum with infinite second moment.

The principal cutoff result separates the extreme-action scale \(\sqrt n\)
from the larger fluctuation scale \(\sqrt{n\log n}\). It identifies the exact
limiting retained fraction, its first prefix-independent correction, and
the first correction to the minimal budget. A finite positive inequality
then supplies a logically independent route to certified truncation.
Finally, the moving finite folds reproduce the same leading action cost
through their critical-value drift, with their own convergent reduced chart.

These results constitute a proved model-specific extension of the inspected
repository program. They are not a claim of established publication priority
or a resolution of unrestricted transseries realization. The explicit
coefficient, support, and error formulas make the next steps---higher
corrections, uniform endpoint crossover, and proof-producing computation---
precise enough to test and formalize.

\phantomsection
\addcontentsline{toc}{section}{References}
\begin{thebibliography}{9}
\bibitem{repo-index}
Vladimir Reshetnikov, \emph{ProveIt: Transseries project and series-and-transseries
index}, repository documentation, observed revision
\texttt{3f2d57344dc6f0b3dfa376f9490bfe13d1beb81e}, 29 September 2026.
\url{https://github.com/VladimirReshetnikov/ProveIt/blob/3f2d57344dc6f0b3dfa376f9490bfe13d1beb81e/Analysis/Transseries/docs/series-and-transseries/README.md}.
Inspection scope is stated in Section~\ref{sec:intro}.

\bibitem{repo-fold}
\emph{Critical Transseries at a Moving Fold: a uniform two-sheet atlas,
large-sector crossover, and resonant action splitting}, research package in
\repo, 29 September 2026. Package description inspected in the group index
\cite{repo-index}; the present article does not claim a full audit of that
manuscript.
\url{https://github.com/VladimirReshetnikov/ProveIt/tree/3f2d57344dc6f0b3dfa376f9490bfe13d1beb81e/Analysis/Transseries/docs/series-and-transseries/Critical_Transseries_Moving_Fold}.

\bibitem{repo-regularity}
\emph{A Sharp Regularity Classification for Countable Exponential-Feedback
Transseries}, research draft prepared for Vladimir Reshetnikov,
29 September 2026. Package README inspected at the revision above.
\url{https://github.com/VladimirReshetnikov/ProveIt/blob/3f2d57344dc6f0b3dfa376f9490bfe13d1beb81e/Analysis/Transseries/docs/series-and-transseries/Exponential_Feedback_Regularity_Classification/README.md}.

\bibitem{prior-stable}
\emph{Beyond Finite-Action Folds: Critical Hahn Transseries, Stable Sector
Asymptotics, and Sharp Action Budgets}, unpublished companion research draft
prepared for Vladimir Reshetnikov, 29 September 2026. Inspected in the user's
file library under the filename \texttt{article(20260929-193550).tex}.
Used only for scope and overlap comparison; no public publication or
independent verification is asserted.

\bibitem{janson}
Svante Janson, Simply generated trees, conditioned Galton--Watson trees,
random allocations and condensation, \emph{Probability Surveys} \textbf{9}
(2012), 103--252. DOI: \texttt{10.1214/11-PS188}.
\url{https://arxiv.org/abs/1112.0510}.

\bibitem{chakrabarty-samorodnitsky}
Arijit Chakrabarty and Gennady Samorodnitsky, Understanding heavy tails in a
bounded world or, is a truncated heavy tail heavy or not?,
\emph{Stochastic Models} \textbf{28}(1) (2012), 109--143.
DOI: \texttt{10.1080/15326349.2012.646551}.
\url{https://arxiv.org/abs/1001.3218}.

\bibitem{dlmf-polylog}
NIST Digital Library of Mathematical Functions, \S25.12,
\emph{Polylogarithms}, consulted 29 September 2026.
\url{https://dlmf.nist.gov/25.12}.
The integer-order logarithmic expansion used here is also derived directly
in Lemma~\ref{lem:normal}.

\bibitem{dlmf-w}
NIST Digital Library of Mathematical Functions, \S4.13,
\emph{Lambert W-Function}, consulted 29 September 2026.
\url{https://dlmf.nist.gov/4.13}.
The local convergence arguments needed here are proved explicitly.
\end{thebibliography}
\end{document}
